\chapter{Vector and tensor fields}
\label{chap:vector-tensor-fields}

\begin{quote}
Spec dir:
\passthrough{\lstinline!specs/2026-05-04-phase-13-vector-tensor-fields/!}.
Version: \passthrough{\lstinline!0.13.0!}.
\end{quote}

\subsection{What this gives you}\label{phase13-what}

\passthrough{\lstinline!Field<T>!} now instantiates with a rank-2
tensor element, and the rank-raising / rank-lowering arithmetic
identities work as expected. Three new public symbols:

\begin{itemize}
\tightlist
\item
  \passthrough{\lstinline!gridcalc::core::Mat3d!} — alias for
  \passthrough{\lstinline!Eigen::Matrix3d!} (standard Eigen module).
  Lives in
  \passthrough{\lstinline!core/EigenAliases.hpp!} alongside
  \passthrough{\lstinline!Vec3d!}.
\item
  \passthrough{\lstinline!diff::gradient(Field<Vec3d>) -> Field<Mat3d>!} —
  rank-raising gradient with the convention
  \(M(i, j) = \partial_j v_i\).
\item
  \passthrough{\lstinline!diff::divergence(Field<Mat3d>) -> Field<Vec3d>!} —
  rank-lowering divergence with the matched convention
  \((\mathrm{div}\,M)_i = \partial_j M(i, j)\).
\end{itemize}

Plus the contraction primitives in the new
\passthrough{\lstinline!tensor/Contraction.hpp!}:

\begin{itemize}
\tightlist
\item
  \passthrough{\lstinline!tensor::trace(Field<Mat3d>) -> Field<double>!}
\item
  \passthrough{\lstinline!tensor::singleContract(Field<Mat3d>, Field<Mat3d>) -> Field<Mat3d>!}
  — pointwise matrix product \((A \cdot B)(i,j) = A(i,k)\,B(k,j)\).
\item
  \passthrough{\lstinline!tensor::doubleContract(Field<Mat3d>, Field<Mat3d>) -> Field<double>!}
  — pointwise full contraction \(A : B = A(i,j)\,B(i,j)\).
\end{itemize}

\textbf{Out of scope at Phase 13.}
\passthrough{\lstinline!Field<Tensor3>!} and
\passthrough{\lstinline!Field<Tensor4>!} (rank-3 and rank-4 tensor
fields) are deferred to Phase 15, where the contraction
expression-template work that motivates them lands together. The
unsupported-Eigen-tensor-module portability question stays deferred.

\subsection{Index conventions}\label{phase13-conventions}

The rank-2 gradient stores
\(M(i, j) = \partial_j v_i\): row index \(i\) selects the vector
component being differentiated; column index \(j\) selects the
derivative axis. Under this convention the trace recovers the
divergence,
\[
\mathrm{tr}(M) = \partial_x v_x + \partial_y v_y + \partial_z v_z = \nabla\cdot v,
\]
and the matched rank-2 divergence
\((\mathrm{div}\,M)_i = \partial_j M(i, j)\) gives the natural vector
Laplacian under composition:
\[
\big(\mathrm{div}(\nabla v)\big)_i = \partial_j \partial_j v_i = (\nabla^2 v)_i.
\]

\subsection{API}\label{phase13-api}

The Phase 2 scalar gradient and divergence keep their existing
signatures unchanged; Phase 13 adds new overloads on
\passthrough{\lstinline!Field<Vec3d>!} and
\passthrough{\lstinline!Field<Mat3d>!} respectively. Both new
overloads inherit the Phase 7
\passthrough{\lstinline!Order!} template parameter (default
\passthrough{\lstinline!2!}; \passthrough{\lstinline!4!}-th order
available):

\begin{lstlisting}[language={C++}]
namespace gridcalc::diff {

template <int Order = 2>
core::Field<core::Mat3d> gradient(const core::Field<core::Vec3d>& field);  // new

template <int Order = 2>
core::Field<core::Vec3d> divergence(const core::Field<core::Mat3d>& mfield);  // new

}  // namespace gridcalc::diff

namespace gridcalc::tensor {

core::Field<double>      trace(const core::Field<core::Mat3d>& mfield);
core::Field<core::Mat3d> singleContract(const core::Field<core::Mat3d>& a,
                                        const core::Field<core::Mat3d>& b);
core::Field<double>      doubleContract(const core::Field<core::Mat3d>& a,
                                        const core::Field<core::Mat3d>& b);

}  // namespace gridcalc::tensor
\end{lstlisting}

\subsection{Worked example: three vector
identities}\label{phase13-identities}

The Phase 13 acceptance test
(\href{../../../test/vector_tensor_test.cpp}{\passthrough{\lstinline!test/vector\_tensor\_test.cpp!}})
verifies three identities programmatically.

\subsubsection{1. Hessian symmetry
(\(\nabla\times\nabla\phi = 0\))}\label{phase13-hessian}

For smooth \(\phi\), the Hessian \(\nabla\nabla\phi\) is symmetric.
Composed Phase 2 gradients realize the rank-raising version directly:

\begin{lstlisting}[language={C++}]
Field<double> phi(grid, [](double x, double y, double z) {
    return std::sin(x) * std::cos(y) * std::sin(z);
});
auto H = gradient(gradient(phi));  // Field<Vec3d> -> Field<Mat3d>

// Each H(i, j, k) is symmetric to round-off:
EXPECT_NEAR(H(7, 11, 5)(0, 1), H(7, 11, 5)(1, 0), 1e-12);
\end{lstlisting}

The agreement is to round-off (not stencil-order tolerance) because
the FD operators \(\partial_i\) and \(\partial_j\) commute on a
periodic uniform grid.

\subsubsection{2.
\(\mathrm{tr}(\nabla v) = \nabla\cdot v\)}\label{phase13-trace}

\begin{lstlisting}[language={C++}]
auto v = sampleVectorField(grid, [](double x, double y, double z) {
    return Vec3d(std::sin(x) * std::cos(y),
                 std::sin(y) * std::cos(z),
                 std::sin(z) * std::cos(x));
});
auto trace_grad_v = tensor::trace(diff::gradient(v));
auto div_v        = diff::divergence(v);
// trace_grad_v == div_v to round-off pointwise.
\end{lstlisting}

\subsubsection{3. \(\nabla\cdot(\nabla v) = \nabla^2
v\)}\label{phase13-vec-laplacian}

For the same \(v\), each Cartesian component depends on exactly two
axes, so its scalar Laplacian sums two
\(-\text{component}\) contributions and the vector Laplacian is
\(\nabla^2 v = -2\,v\). The composition
\passthrough{\lstinline!divergence(gradient(v))!} is a stride-2
second-derivative stencil (the same composed-stencil phenomenon
documented for \passthrough{\lstinline!divergence(gradient(\(\psi\)))!} versus
\passthrough{\lstinline!laplacian(\(\psi\))!} in Phase 2's note); it converges
to the analytical \(\nabla^2 v\) at \(\mathcal{O}(h^2)\) but with a
different leading-error constant than
\passthrough{\lstinline!diff::laplacian!} applied componentwise. The
test checks the rate against the closed-form
\(-2\,v\), not against a discrete componentwise Laplacian.

\subsection{What this does \emph{not} do}\label{phase13-not-do}

\begin{itemize}
\tightlist
\item
  \textbf{Rank-3 / rank-4 tensor fields.} Deferred to Phase 15.
\item
  \textbf{Expression-template contraction layer.} Deferred to Phase 15
  alongside the
  \passthrough{\lstinline!func::evaluate!}-over-tensor-fields surface
  that motivates it.
\item
  \textbf{Explicit \texttt{curl} operator.} The
  \passthrough{\lstinline!gradient(gradient(phi))!}-symmetry test
  discharges the
  \(\nabla\times\nabla\phi = 0\) identity in rank-raising
  vocabulary; an explicit
  \passthrough{\lstinline!curl!} (rank-1 \(\to\) rank-1) is not on
  the roadmap and is not added speculatively.
\item
  \textbf{Vector / tensor spectral cross-check.} The Phase 9 fixture
  parameterizes over scalar-input FD operators; the rank-raising
  variants are not (yet) in its scope.
\end{itemize}

For the construction of the manufactured solutions and the rate
arguments, see the ``Vector and tensor fields'' chapter of the
Developer Note.
